\chapter{Conclusiones, competencias y trabajo futuro}
\label{cap:conclusiones}

El Capítulo~\ref{cap:introduccion} planteó el problema de partida, la
motivación del trabajo y un objetivo general acompañado de diez objetivos
específicos. Los capítulos intermedios han desarrollado, capa a capa, la
arquitectura Open CVN que responde a ese planteamiento, y el
Capítulo~\ref{cap:evaluacion} ha aportado evidencia reproducible de su
comportamiento. El presente capítulo cierra el documento retomando esos
objetivos para valorar su grado de cumplimiento, relacionando el trabajo
realizado con las competencias de Computación implicadas, y señalando las
limitaciones y líneas de trabajo futuro que se derivan de todo lo anterior.

\section{Resumen del trabajo}
\label{sec:conclusiones_resumen}

\textbf{Problema y motivación.} CVN proporciona a España una norma
consolidada para representar currículos académicos e investigadores, pero
su uso automatizado sigue siendo costoso: su representación XML oficial y
la herramienta asociada no se diseñaron como base directa de aplicaciones
abiertas y ligeras, tal y como se planteó en el
Capítulo~\ref{cap:introduccion} y se analizó con evidencia técnica concreta
en el Capítulo~\ref{cap:ecosistema_cvn}.

\textbf{Arquitectura desarrollada.} Frente a ese problema, este trabajo
construye Open CVN: una arquitectura por capas, presentada en el
Capítulo~\ref{cap:metodologia}, que separa el paquete oficial CVN, los
metadatos normalizados, la política semántica, los modelos de dominio, el
modelo conceptual y el formato Open CVN JSON. El
Capítulo~\ref{cap:pipeline} implementa la generación estructural, la
normalización y la política semántica que alimentan esa arquitectura, y el
Capítulo~\ref{cap:formato_herramienta} completa el recorrido con el
formato Open CVN JSON, su parser, su validador y una herramienta local que
demuestra su uso.

\textbf{Resultados.} El Capítulo~\ref{cap:evaluacion} evalúa esta
arquitectura con una batería de 488 pruebas automatizadas y reproducibles,
organizada en ocho niveles alineados con la propia arquitectura, sin
ningún fallo en su última ejecución. Los resultados confirman garantías
fuertes allí donde el proceso es determinista (generación estructural,
normalización, política semántica y validación del formato Open CVN JSON)
y garantías parciales, documentadas de forma explícita, donde el propio
dominio o el alcance del trabajo lo exigen (importación de CVN XML e
importación asistida por LLM).

\textbf{Visión original y su alcance.} La motivación de este trabajo,
presentada en el Capítulo~\ref{cap:introduccion}, se planteó originalmente
en cuatro pasos orientados al uso final. La Tabla~\ref{tab:conclusiones_vision}
retoma cada uno de ellos y valora en qué medida Open CVN se ha acercado a
él.

\begin{table}[htbp]
  \centering
  \footnotesize
  \renewcommand{\arraystretch}{1.45}
  \caption{Visión original del proyecto y grado de alcance logrado por Open CVN.}
  \label{tab:conclusiones_vision}
  \begin{tabular}{p{0.34\textwidth} p{0.54\textwidth}}
    \hline
    \textbf{Paso de la visión original} & \textbf{Alcance logrado} \\
    \hline
    Disponer de una especificación de currículo abierta &
    Cumplido: el formato Open CVN JSON y su JSON Schema definen esa
    especificación de forma validable \\
    Poder construir almacenes de currículos abiertos, con una forma
    definida de leer y escribir en ellos &
    Cumplido parcialmente: el parser, el validador y el almacenamiento
    local de la herramienta de gestión lo demuestran, pero como una
    aplicación de un único usuario, no como una infraestructura de
    almacenes abiertos compartidos entre aplicaciones \\
    Poder generar, a partir de un currículo en esa especificación,
    currículos adaptados a necesidades concretas &
    Cumplido: mediante la exportación a LaTeX y PDF y la selección de
    secciones y entradas entre el currículo maestro y sus versiones
    derivadas \\
    Poder importar a la especificación currículos ya existentes &
    Cumplido con garantías parciales y explícitamente documentadas:
    importación determinista desde CVN XML y CVN PDF, con recurso
    opcional a modelos de lenguaje de gran tamaño \\
    \hline
  \end{tabular}
\end{table}

\section{Cumplimiento de los objetivos}
\label{sec:conclusiones_objetivos}

El objetivo general planteado en el Capítulo~\ref{cap:introduccion} consistía
en diseñar e implementar una arquitectura computacional abierta capaz de
representar, validar, transformar, almacenar y exportar currículos
académicos e investigadores a partir de CVN, sin que la representación XML
oficial ni la herramienta oficial condicionaran por completo el modelo
interno del sistema. Este objetivo se considera cumplido: la arquitectura
por capas presentada en el Capítulo~\ref{cap:metodologia} e implementada en
los Capítulos~\ref{cap:pipeline} y~\ref{cap:formato_herramienta} separa de
forma explícita el paquete oficial CVN, los metadatos normalizados, el
modelo conceptual y el formato Open CVN JSON, y la evaluación del
Capítulo~\ref{cap:evaluacion} confirma que esa separación funciona de forma
reproducible.

Únicamente uno de los objetivos específicos se cumple con garantías
parciales, y de forma deliberada: la importación desde CVN XML es, por
diseño, un mapeo semántico parcial y creciente, y la importación asistida
por LLM es opcional y no se considera autoritativa sin revisión, tal y como
se discutió con evidencia en el Capítulo~\ref{cap:evaluacion}. El resto de
objetivos se apoya en mecanismos deterministas que la batería de pruebas
verifica de forma reproducible. La Tabla~\ref{tab:conclusiones_objetivos}
detalla el grado de cumplimiento de cada objetivo específico, en el mismo
orden en que se presentaron en el Capítulo~\ref{cap:introduccion}.

\begin{table}[htbp]
  \centering
  \footnotesize
  \renewcommand{\arraystretch}{1.5}
  \caption{Grado de cumplimiento de los objetivos específicos.}
  \label{tab:conclusiones_objetivos}
  \begin{tabular}{p{0.06\textwidth} p{0.52\textwidth} p{0.3\textwidth}}
    \hline
    \textbf{Obj.} & \textbf{Descripción resumida} & \textbf{Grado de cumplimiento} \\
    \hline
    OE1 & Análisis del ecosistema CVN y sus artefactos oficiales & Cumplido \\
    OE2 & Estudio de alternativas de representación, modelado y validación & Cumplido \\
    OE3 & Generación estructural reproducible desde XSD & Cumplido \\
    OE4 & Normalización de metadatos CVN & Cumplido \\
    OE5 & Definición de reglas semánticas trazables & Cumplido \\
    OE6 & Generación de modelos de dominio y artefactos conceptuales & Cumplido \\
    OE7 & Definición del formato Open CVN JSON & Cumplido \\
    OE8 & Parser, validador, almacenamiento y exportación & Cumplido \\
    OE9 & Importación determinista y asistida por LLM & Cumplido con garantías parciales \\
    OE10 & Verificación automatizada y flujos reproducibles & Cumplido \\
    \hline
  \end{tabular}
\end{table}

\section{Contribuciones principales}
\label{sec:conclusiones_contribuciones}

Las contribuciones anunciadas en el Capítulo~\ref{cap:introduccion} quedan
demostradas, con evidencia concreta, a lo largo del documento.

La primera es el \textbf{análisis computacional del ecosistema CVN}: el
Capítulo~\ref{cap:ecosistema_cvn} documenta con evidencia concreta las
cuatro capas del paquete oficial y sus tres familias auxiliares, junto con
las limitaciones técnicas detectadas, en lugar de asumir que dicho paquete
puede tratarse como una fuente homogénea y sin fricciones.

La segunda es la \textbf{arquitectura reproducible por capas} presentada en
el Capítulo~\ref{cap:metodologia} e implementada en el
Capítulo~\ref{cap:pipeline}: la separación entre bindings generados,
normalización, política semántica, modelos de dominio, modelo conceptual,
formato JSON y herramienta local evita que una única representación
concentre todas las responsabilidades del sistema.

La tercera es la \textbf{trazabilidad entre CVN y Open CVN}: el
Capítulo~\ref{cap:pipeline} presenta un mecanismo concreto de trazabilidad,
con anotaciones no validantes y un recorrido completo y verificable desde
un código CVN hasta su definición final en el formato de salida.

La cuarta es el propio \textbf{formato Open CVN JSON}, cuya estructura raíz,
organización por áreas conceptuales y forma común de las referencias
controladas se definen en el Capítulo~\ref{cap:formato_herramienta}.

La quinta son las \textbf{herramientas de validación y uso}: el parser, el
validador y la herramienta local de gestión, con su modelo de currículo
maestro y versiones derivadas, descritos en el mismo
Capítulo~\ref{cap:formato_herramienta}.

La sexta es la \textbf{importación determinista y asistida por LLM}: las
vías de importación desde CVN XML y CVN PDF, con el fallback opcional por
LLM sujeto a validación posterior, presentadas también en el
Capítulo~\ref{cap:formato_herramienta}.

La séptima y última es la \textbf{verificación automatizada}: la batería
de 488 pruebas automatizadas, organizada en ocho niveles alineados con la
arquitectura, documentada con resultados reproducibles en el
Capítulo~\ref{cap:evaluacion}.

Estas contribuciones no sustituyen al ecosistema CVN oficial, sino que
proponen una capa complementaria que facilita su tratamiento computacional
en contextos donde son necesarias la validación formal, la transformación
entre formatos y la gestión automatizada de currículos.

\section{Competencias desarrolladas}
\label{sec:conclusiones_competencias}

El Capítulo~\ref{cap:introduccion} identificó cuatro competencias de la
intensificación de Computación trabajadas en este proyecto y anunció que su
grado de cumplimiento se retomaría en las conclusiones. La
Tabla~\ref{tab:conclusiones_competencias} recoge esa evidencia concreta.

\begin{table}[htbp]
  \centering
  \footnotesize
  \renewcommand{\arraystretch}{1.55}
  \caption{Competencias de Computación y evidencia concreta de cada una.}
  \label{tab:conclusiones_competencias}
  \begin{tabular}{p{0.09\textwidth} p{0.83\textwidth}}
    \hline
    \textbf{Comp.} & \textbf{Evidencia concreta} \\
    \hline
    CM1 & El diseño de la arquitectura por capas del
    Capítulo~\ref{cap:metodologia} y la definición del formato Open CVN JSON
    del Capítulo~\ref{cap:formato_herramienta} como artefactos técnicos
    nuevos derivados de un análisis conceptual del dominio curricular. \\
    CM2 & El procesamiento léxico, sintáctico y semántico de los esquemas y
    documentos CVN a lo largo del pipeline del
    Capítulo~\ref{cap:pipeline}: generación estructural desde XSD,
    normalización de metadatos y política semántica que interpreta
    y traduce dicha estructura. \\
    CM5 & La formalización del conocimiento curricular contenido en el
    paquete oficial CVN en un modelo conceptual computable y validable,
    desarrollada en los Capítulos~\ref{cap:ecosistema_cvn}
    a~\ref{cap:formato_herramienta}. \\
    CM6 & La herramienta local de gestión del
    Capítulo~\ref{cap:formato_herramienta}, que presenta información
    curricular compleja y permite seleccionar, versionar y exportar
    secciones concretas de un currículo mediante una interfaz de línea de
    órdenes. \\
    \hline
  \end{tabular}
\end{table}

\section{Limitaciones y trabajo futuro}
\label{sec:conclusiones_limitaciones}

Los capítulos anteriores han documentado limitaciones concretas en distintos
puntos del sistema, sin ocultarlas ni presentarlas como defectos de
implementación cuando, en realidad, proceden del propio paquete oficial o de
una decisión deliberada de alcance, tal y como ya se distinguió en el
Capítulo~\ref{cap:evaluacion}. Esta sección retoma cada una de ellas con el
detalle suficiente para justificar la línea de trabajo futuro que le
corresponde, antes de consolidarlas en una tabla de referencia rápida.

La primera limitación son las \textbf{inconsistencias entre artefactos del
paquete oficial CVN}, como la discrepancia puntual entre el XML canónico del
modelo en árbol y su esquema de validación, o la deriva de empaquetado entre
las familias auxiliares descrita en el Capítulo~\ref{cap:ecosistema_cvn}.
Esta limitación no es atribuible al sistema desarrollado, sino al propio
material fuente, por lo que el trabajo futuro correspondiente no puede
consistir en corregirla, sino en revisarla si el paquete oficial se
actualiza y documentar cualquier discrepancia nueva con el mismo criterio de
evidencia ya aplicado.

La segunda es la \textbf{expresividad limitada de algunos bindings
estructurales generados}: construcciones como el mecanismo \textit{choice}
de XSD, que no siempre se traduce en una exclusión mutua real, o los
atributos generados con un tipo demasiado genérico, tal y como se analizó en
el Capítulo~\ref{cap:ecosistema_cvn}. Estas limitaciones se confinan de
forma deliberada a la capa estructural generada y no se propagan al formato
Open CVN JSON, pero el trabajo futuro puede ampliar las validaciones de
dominio que compensan esta pérdida de expresividad fuera de dicha capa.

La tercera son las \textbf{referencias auxiliares no resueltas con el
paquete actual}, de las que el caso de \texttt{CVN\_AGENCY\_C} es el
ejemplo verificado en el Capítulo~\ref{cap:pipeline}: una tabla citada por
el manual sin equivalente localizable en el resto del paquete. El sistema
conserva esta referencia como no resuelta en lugar de forzar una
interpretación no justificada, y el trabajo futuro solo puede mejorar la
cobertura de tablas auxiliares si aparece evidencia oficial adicional que
hoy no existe.

La cuarta es la \textbf{importación de CVN XML semánticamente parcial}: un
mapeo creciente sobre los casos ya cubiertos, y no un conversor completo de
cualquier documento CVN XML real, según se discutió con evidencia en el
Capítulo~\ref{cap:evaluacion}. A diferencia de las dos limitaciones
anteriores, esta sí admite una línea de trabajo futuro directa y
sustancial: completar el mapeo semántico sobre más casos reales y evaluar
el sistema con currículos reales, en lugar de únicamente con los ejemplos
sintéticos empleados en este trabajo.

La quinta es que el \textbf{JSON Schema no puede expresar invariantes entre
atributos}, porque un esquema de validación declarativo, igual que un
diagrama de clases UML, solo puede describir la estructura, el tipo y la
cardinalidad de un atributo aislado, no reglas que combinen varios atributos
entre sí. El modelo conceptual actual ya contiene, sin formalizar, patrones
de este tipo: un campo de valor controlado seguido de un campo de texto
libre para el caso ``otros'', donde la regla de negocio implícita exige que
ese texto libre sea obligatorio únicamente cuando se elige dicha opción; o
un par de fechas de inicio y de fin dentro de una misma entidad, donde se
espera que la primera preceda a la segunda. Ninguna de estas reglas está
implementada actualmente ni en los modelos generados ni en los diagramas
conceptuales, por lo que el trabajo futuro correspondiente, estudiar la
incorporación de restricciones OCL sobre el modelo conceptual, constituiría
una extensión deliberada de dicho modelo, y no la transcripción de una
validación ya existente.

La sexta y última es la \textbf{generación de PDF dependiente de un motor
LaTeX externo}, ya señalada en el Capítulo~\ref{cap:formato_herramienta}:
la exportación aprovecha una cadena de herramientas LaTeX ya madura en
lugar de reimplementar un renderizador de PDF propio, a costa de exigir
dicho motor en el entorno de ejecución. El trabajo futuro puede explorar
mecanismos de generación adicionales que reduzcan esta dependencia sin
renunciar a la calidad tipográfica de la exportación actual.

Además de estas seis líneas, ligadas cada una a una limitación concreta, el
trabajo admite otras extensiones de alcance no derivadas de una limitación
específica, sino de una ampliación natural del sistema desarrollado:
desarrollar una interfaz gráfica si se considera necesaria más allá de la
herramienta de línea de órdenes actual, estudiar extensiones semánticas
como JSON-LD para facilitar la integración con datos enlazados, e
investigar la integración con fuentes institucionales de información
curricular.

La Tabla~\ref{tab:conclusiones_limitaciones} consolida, a modo de
referencia rápida, las seis limitaciones ya explicadas junto con la línea
de trabajo futuro que se deriva de cada una.

\begin{table}[htbp]
  \centering
  \footnotesize
  \renewcommand{\arraystretch}{1.55}
  \caption{Limitaciones documentadas y trabajo futuro asociado.}
  \label{tab:conclusiones_limitaciones}
  \begin{tabular}{p{0.43\textwidth} p{0.43\textwidth}}
    \hline
    \textbf{Limitación} & \textbf{Trabajo futuro asociado} \\
    \hline
    Inconsistencias entre artefactos del paquete oficial CVN (discrepancia
    XML/XSD, deriva de empaquetado) & Revisar estas inconsistencias si el
    paquete oficial se actualiza, documentando cualquier discrepancia nueva
    con el mismo criterio de evidencia \\
    Expresividad limitada de algunos bindings estructurales generados
    (\textit{choice} sin exclusión mutua, tipos genéricos) & Ampliar las
    validaciones de dominio fuera de la capa estructural generada \\
    Referencias auxiliares no resueltas con el paquete actual & Mejorar la
    cobertura de tablas auxiliares si aparece evidencia oficial adicional \\
    Importación de CVN XML semánticamente parcial & Completar el mapeo
    semántico sobre más casos reales y evaluar el sistema con currículos
    reales \\
    JSON Schema sin capacidad de expresar invariantes entre atributos &
    Estudiar la incorporación de restricciones OCL sobre el modelo
    conceptual \\
    Generación de PDF dependiente de un motor LaTeX externo & Explorar
    mecanismos de generación adicionales que reduzcan esta dependencia \\
    \hline
  \end{tabular}
\end{table}

\section{Conclusiones}
\label{sec:conclusiones_finales}

El resumen, el cumplimiento de objetivos y las contribuciones ya
presentados en este capítulo permiten cerrar el documento con una
valoración conjunta, en lugar de con la simple repetición de lo ya
expuesto. El objetivo general se considera cumplido y, con él, los diez
objetivos
específicos que lo desarrollaban, con la única salvedad deliberada de la
importación asistida por LLM, cuyo carácter opcional y no autoritativo es
una decisión de diseño y no una carencia. Esta conclusión no se apoya en
una afirmación sin respaldo: el Capítulo~\ref{cap:evaluacion} la sostiene
con una batería de 488 pruebas automatizadas y reproducibles, organizada
según las mismas capas que estructuran el documento, lo que permite
sostener que las garantías descritas no son solo argumentativas, sino
verificables por cualquier persona que ejecute el mismo comando sobre el
mismo repositorio.

Las limitaciones recogidas en este capítulo no comprometen esta conclusión:
delimitan, con precisión y sin ocultamiento, hasta dónde llega la propuesta
actual y qué preguntas quedan abiertas para su continuación. Algunas de
ellas son inherentes al propio paquete oficial CVN y quedan fuera del
control del sistema desarrollado; otras son decisiones deliberadas de
alcance, propias de un Trabajo Fin de Grado y no de un producto cerrado. En
ambos casos, el sistema las documenta en lugar de disimularlas, lo que
constituye en sí mismo una forma de rigor técnico.

En conjunto, este trabajo demuestra que es posible construir, a partir del
ecosistema CVN, una arquitectura computacional abierta que separa la
interoperabilidad estructural del significado curricular, conserva
trazabilidad verificable hacia sus fuentes oficiales y ofrece garantías
reproducibles allí donde el proceso es determinista. Esa arquitectura, más
que el formato JSON o la herramienta local que la demuestran, es la
aportación principal de este Trabajo Fin de Grado a la intensificación de
Computación en la que se enmarca.
